\scanpage{12}
\[
\sigma_\infty(g^{-1})=l_0^\beta g l_1^\beta.
\]
Ce qui détermine $\beta$ de façon unique en fonction de $g$.

\deleted{Appliquant la relation fondamentale
\[
u\rho=\operatorname{int}(g)\rho u
\]
à $L_0$, on trouve, puisque $u(L_0)=L_0$, $\rho(L_0)=L_1$
\[
u(L_1)=\operatorname{int}(g)(L_{\unsure{1|0}})
\]
et ceci détermine déjà \deleted{\unsure{$\operatorname{int}(g)$}} modulo multiplication à droite par un él\textsuperscript{t} de $L_1$, \unsure{par} la donnée de $u$.}

On a des vues plus géométriques \unsure{de} \inserted{\unsure{sur}} l'extension de $\unsure{\mathfrak S_3|\mathcal S_3}$ par $\pi$ \inserted{dans $\operatorname{Aut}_{\mathrm{lac}}(\pi)$,} que les six générateurs $l_0,l_1,l_\infty,\sigma_0,\sigma_1,\sigma_\infty$ (donnant $\rho$ par $\sigma_0\sigma_1=l_1^{-1}\rho$ \deleted{\unsure{p.ex.}} i.e. $\rho=\sigma_0\sigma_1l_1$ p.ex.), en introduisant les groupes de ramification dans les pts $0,1,\infty$, qui sont canoniq\textsuperscript{t} isomorphes à $\widehat{\mathbb Z}$.
On définira donc $\varepsilon_0,\varepsilon_1,\varepsilon_\infty$ par les relations
\[
\varepsilon_0^2=l_0,\qquad\varepsilon_1^2=l_1,\qquad\varepsilon_\infty^2=l_\infty
\]
et plus précisément, $\varepsilon_i$ sera l'unique \deleted{\unsure{él\textsuperscript{t}}} \inserted{autom.} dans la classe mod $\pi$ des $\sigma_i$, \deleted{et} qui fixe $l_i$ et est \deleted{\unsure{de|$\sigma$}} \inserted{\unsure{de carré $l_i$}}
\mspage{11}
